% year: תשפ"ג
% semester: א
% moed: ג
% source: exams/מבחנים/תשפג/מועד ג תשפג.pdf
% number: 3ב
% categories: func-analysis

\begin{question}
מצאו תחומי קמירות ונקודות פיתול עבור הפונקציה $f(x)=e^{-x^2+2x}$
\end{question}

\begin{solution}
$f$ מוגדרת וגזירה פעמיים על כל $\R$. לפי כלל השרשרת
\[ f'(x)=(2-2x)e^{-x^2+2x}, \]
\[ f''(x)=-2e^{-x^2+2x}+(2-2x)^2e^{-x^2+2x}=(4x^2-8x+2)e^{-x^2+2x}=2(2x^2-4x+1)e^{-x^2+2x}. \]
מכיוון ש-$e^{-x^2+2x}>0$, סימן $f''$ הוא סימן $2x^2-4x+1$, שאפסיו
\[ x_{1,2}=\frac{4\pm\sqrt{16-8}}{4}=1\pm\frac{\sqrt2}{2}. \]
זו פרבולה עם מקדם מוביל חיובי, ולכן:
\begin{itemize}
  \item $f''>0$ ($f$ \textbf{קמורה}) עבור $x<1-\frac{\sqrt2}{2}$ ועבור $x>1+\frac{\sqrt2}{2}$;
  \item $f''<0$ ($f$ \textbf{קעורה}) עבור $1-\frac{\sqrt2}{2}<x<1+\frac{\sqrt2}{2}$.
\end{itemize}
בשתי הנקודות $x=1\pm\frac{\sqrt2}{2}$ הנגזרת השנייה מחליפה סימן והפונקציה רציפה, ולכן אלה נקודות פיתול. מכיוון ש-$-x^2+2x=1-(x-1)^2=1-\frac12=\frac12$ בנקודות אלה, נקודות הפיתול הן
\[ \boxed{\left(1-\tfrac{\sqrt2}{2},\ \sqrt e\right),\quad \left(1+\tfrac{\sqrt2}{2},\ \sqrt e\right)}. \]
\end{solution}
